\section{Results and Discussion}
To determine the most viable approach for addressing the ICT infrastructure deficit in the Municipality of Maigo, this study evaluates three distinct policy alternatives. These alternatives are assessed based on three criteria: Economic Feasibility (Capital Expenditure vs. Operating Expenditure), Technical Feasibility (Maintenance and longevity), and Socio-Economic Impact (Value for Money and MSME resilience).

\subsection{Alternative A: Do nothing (Status Quo)}
Under this alternative, the Local Government Unit (LGU) and educational institutions maintain their current operations. MSU-MCEST and TESDA continue to rely on existing, limited laboratory space, while out-of-school youth remain without access to advanced ICT certifications.

\begin{itemize}
    \item Economic Impact: While this requires zero government spending initially, the opportunity cost is severe. The municipality loses potential skilled labor.
    \item Socio-Economic Impact: Local MSMEs (internet cafés) will continue to face declining revenues and potential bankruptcy due to underutilized hardware during standard business hours.
    \item Verdict: Highly unsustainable and fails to address the digital divide.
\end{itemize}

\subsection{Alternative B: LGU builds new public labs (High CapEx)}
This is the standard approach where the LGU or national government funds the construction of new, dedicated public computer laboratories.

\begin{itemize}
    \item Economic Impact: This requires massive Capital Expenditure (CapEx). A standard 30-seat laboratory requires funding for high-end PCs, air conditioning, structural security, and commercial-grade networking equipment. Furthermore, IT hardware depreciates rapidly, requiring hardware replacement cycles every 3 to 5 years.
    \item Technical Impact: Requires the LGU or schools to hire dedicated, full-time IT administrators for continuous maintenance.
    \item Verdict: While it solves the educational deficit, it is economically prohibitive for a standard municipality and does nothing to support the local MSME sector.
\end{itemize}

\subsection{Alternative C: The Dual-Use PPP Model (Zero-CapEx)}
This alternative represents the proposed framework: institutionalizing a service-contracting model under the PPP Code (RA 11966) where the LGU leases existing private cybercafés during the day.

\begin{itemize}
    \item Economic Impact: Operates on a Zero-CapEx model. The government only incurs Operating Expenditures (OpEx) through hourly or monthly leasing. The financial risk of hardware depreciation and maintenance remains with the private sector.
    \item Socio-Economic Impact: Provides a guaranteed daytime revenue stream for struggling MSMEs, preventing business closures and stabilizing the local economy.
    \item Technical Impact: Utilizing Diskless (PXE-Boot) Multiple Image architecture ensures a sterile, distraction-free educational environment during the day without permanently altering the café's commercial capabilities at night.
    \item Verdict: Highly feasible. It leverages existing, idle community assets to deliver public services efficiently.
\end{itemize}

\subsection{4.4 Comparative Analysis Matrix}
The comparison heavily favors Alternative C as the most viable public policy intervention. By shifting from a "procurement of goods" mindset to a "procurement of services" mindset, the LGU achieves maximum Value for Money (VFM). Alternative C is the only option that simultaneously solves the educational infrastructure deficit while acting as an economic stimulus for local businesses.